\chapter{Verma 模}
\label{ch:05.5}

\section{最高权模\index{zuigaoquanmo@最高权模|hyperpage}范畴}
\label{sec:05.5.1}

固定 $\g$ 中的一个 Borel 子代数 $\bb_0$，并令 $\n_0=[\bb_0,\bb_0]$。设 $\h_0$ 是 $\g$ 中含于 $\bb_0$ 的 Cartan 子代数。根系 $\Sigma$ 经此特化为 $\h_0^*$ 中的根系 $R_0$，而 $R_0^+$ 中正根对应的根子空间张成 $\n_0$。为简化下文的记号，在不致混淆时，我们经此特化把 Cartan 三元组 $(\h,\Sigma,\Sigma^+)$ 与 $(\h_0,R_0,R_0^+)$ 等同起来。记 $\bar{\n}_0$ 为 $R_0$ 中负根对应的根子空间张成的幂零子代数。

设 $V$ 是 $\g$-模。若
\begin{enumerate}[label=(\roman*)]
\item $V$ 有限生成；
\item $V$ 是 $U(\bb_0)$-有限的，即对任意 $v\in V$，$U(\bb_0)v$ 都有限维；
\end{enumerate}
则称 $V$ 是（关于 $\bb_0$ 的）\emph{最高权模}。我们把 $\mathcal{M}_{fg}(U(\g))$ 中由最高权模构成的全子范畴，称为\emph{最高权模范畴}。

设 $V$ 是最高权模。对 $\lambda\in\h_0^*$，置
$$V^\lambda=\{v\in V\mid \text{存在 } k\in\N,\ \text{使对一切 } \xi\in\h_0 \text{ 有 } (\xi-\lambda(\xi))^k v=0\}.$$
则 $V^\lambda$ 是 $V$ 的 $\h_0$-子模，且 $V$ 是诸 $V^\lambda$（$\lambda\in\h_0^*$）的直和。若 $V^\lambda\neq 0$，则称 $\lambda$ 是 $V$ 的\emph{权}。

\begin{lemma}\label{lem:05.1.1}
设 $V$ 是有限生成的 $\g$-模。则下列条件等价：
\begin{enumerate}[label=(\roman*)]
\item $V$ 是最高权模；
\item $V$ 满足：
\begin{enumerate}[label=(\alph*)]
\item $V=\bigoplus_{\lambda\in\h_0^*}V^\lambda$，且各 $V^\lambda$（$\lambda\in\h_0^*$）都有限维；
\item 存在 $V$ 的有限权集 $S_0$，它具有下述性质：对 $V$ 的每个权 $\nu$，都有 $\mu\in S_0$ 使 $\mu-\nu$ 是 $R_0^+$ 中若干根之和。
\end{enumerate}
\end{enumerate}
\end{lemma}

\begin{proof}
设 $V$ 是最高权模。按定义，作为 $\g$-模，$V$ 由某个有限维 $\bb_0$-不变子空间 $U$ 生成。于是由 Poincaré--Birkhoff--Witt 定理，自然映射 $U(\bar{\n}_0)\otimes_\C U\to V$ 是 $\h_0$-模的满态射。这立刻给出 (a) 与 (b)。

现设 $V$ 满足 (a) 与 (b)。取 $v\in V^\lambda$。则 $U(\bb_0)v$ 含于那些 $V^\nu$ 的直和中，其中权 $\nu$ 满足「$\nu-\lambda$ 是若干正根之和」。由 (b)，这样的权只有有限多个。于是由 (a)，$U(\bb_0)v$ 有限维。
\end{proof}

\begin{lemma}\label{lem:05.1.2}
设
$$0\longrightarrow V\longrightarrow V'\longrightarrow V''\longrightarrow 0$$
是 $\g$-模的正合列。则 $V'$ 是最高权模，当且仅当 $V$ 与 $V''$ 都是最高权模。
\end{lemma}

\begin{proof}
显然，若 $V'$ 是最高权模，则 $V$ 与 $V''$ 也都是最高权模。现设 $V$ 与 $V''$ 都是最高权模。则 $V'$ 显然有限生成。又它满足 1.1.(ii) 的条件。故 $V'$ 是最高权模。
\end{proof}

称 $v\in V$ 是 $Z(\g)$-\emph{有限}的，若 $Z(\g)v$ 有限维。显然各 $V^\lambda$ 都是 $Z(\g)$-不变的；又由 1.1.(ii)，其中的向量都是 $Z(\g)$-有限的。这说明 $V$ 中所有向量都是 $Z(\g)$-有限的，即 $V$ 是 $Z(\g)$-有限模。最后，由于 $V$ 有限生成，得到下述结果。

\begin{lemma}\label{lem:05.1.3}
设 $V$ 是最高权模。则 $V$ 在 $Z(\g)$ 中的零化子有限余维。
\end{lemma}

反之，也有下述结论。

\begin{proposition}\label{prop:05.1.4}
设 $\g$-模 $V$ 满足：
\begin{enumerate}[label=(\roman*)]
\item $V$ 有限生成；
\item 对任意 $v\in V$，存在 $k\in\N$ 使 $\n_0^k\cdot v=0$；
\item $V$ 在 $Z(\g)$ 中的零化子有限余维。
\end{enumerate}
则 $V$ 是最高权模。
\end{proposition}

\begin{proof}
设 $U$ 是生成 $V$ 的有限维 $\n_0$-不变子空间。我们对 $\dim U$ 作归纳来证明 $V$ 是最高权模。若 $\dim U=1$，则 $U$ 被 $U(\g)\n_0$ 零化。另一方面，由 Harish-Chandra 同态的性质可知：$Z(\g)\subset U(\h_0)\oplus U(\g)\n_0$ 沿 $U(\g)\n_0$ 到 $U(\h_0)$ 的投影是代数同态，且 $U(\h_0)$ 在其像上有限生成。这立刻推出 $U(\h_0)U=U(\bb_0)U$ 是 $V$ 的有限维子空间。容易验证：由一切使 $U(\bb_0)u$ 有限维的 $u\in V$ 构成的线性子空间 $V'$ 是 $V$ 的 $\g$-子模。由上述讨论它包含 $U$，从而等于 $V$，即 $V$ 是最高权模。现设 $\dim U>1$。由 Engel 定理，$U$ 有一维子空间 $U_0$ 使 $\n_0U_0=0$。设 $V_0$ 是 $V$ 中由 $U_0$ 生成的 $\g$-子模。则由前一段，$V_0$ 是最高权模。令 $V_1=V/V_0$。则 $V_1$ 由 $U_1=U/(U\cap V_0)$ 生成，且 $\dim U_1\le\dim U-1$。故由归纳假设，$V_1$ 是最高权模。再由 1.2 即知 $V$ 是最高权模。
\end{proof}

设 $N_0$ 是 $\operatorname{Int}(\g)$ 中对应于 $\n_0$ 的幂幺子群。则由 1.4，可把 $\n_0$ 的作用指数化为 $N_0$ 的代数作用，从而把最高权模看作 $\mathcal{M}_{fg}(\g,N_0)$ 中的对象。事实上，这样就得到最高权模范畴与 $\mathcal{M}_{fg}(\g,N_0)$ 中「被 $Z(\g)$ 的某个有限余维理想零化」的模构成的全子范畴之间的等同。

下面描述 $\mathcal{M}_{fg}(\g,N_0)$ 中的不可约对象，即不可约最高权模。先构造几个与之密切相关的模。设 $C_\lambda$ 是由 $\lambda\in\h_0^*$ 定义的一维 $\bb_0$-模。若通过右乘把 $U(\g)$ 看作右 $U(\bb_0)$-模，则张量积 $U(\g)\otimes_{U(\bb_0)}C_\lambda$ 通过在第一个因子上的左乘具有自然的左 $U(\g)$-模结构。它显然是最高权模。今置
$$M(\lambda)=U(\g)\otimes_{U(\bb_0)}C_{\lambda-\rho}.$$
最高权模 $M(\lambda)$ 称为由 $\lambda$ 决定的 \emph{Verma 模}。

\begin{lemma}\label{lem:05.1.5}
设 $\lambda\in\h_0^*$。则
\begin{enumerate}[label=(\roman*)]
\item $M(\lambda)$ 的所有权都形如 $\lambda-\rho-\nu$，其中 $\nu$ 是若干正根之和；
\item $\dim M(\lambda)^{\lambda-\rho}=1$；
\item $M(\lambda)$ 有唯一的极大 $\g$-子模 $N(\lambda)$；
\item $N(\lambda)^{\lambda-\rho}=0$。
\end{enumerate}
\end{lemma}

\begin{proof}
由 Poincaré--Birkhoff--Witt 定理，看作 $\h_0$-模时有 $M(\lambda)\cong U(\bar{\n}_0)\otimes_\C C_{\lambda-\rho}$。这立刻给出 (i) 与 (ii)。按 $M(\lambda)$ 的定义，一维子空间 $M(\lambda)^{\lambda-\rho}$ 生成 $M(\lambda)$（作为 $\g$-模）。因此，$M(\lambda)$ 的任何真 $\g$-子模都不含 $M(\lambda)^{\lambda-\rho}$。设 $M$ 是 $M(\lambda)$ 的极大 $\g$-子模，$N$ 是任一异于 $M(\lambda)$ 的 $\g$-子模。则或者 $N\subset M$，或者 $M+N=M(\lambda)$。后一情形会给出
$$M(\lambda)^{\lambda-\rho}=(M+N)^{\lambda-\rho}=M^{\lambda-\rho}+N^{\lambda-\rho}=0,$$
显然不可能。故 $M$ 是唯一的极大 $\g$-子模。
\end{proof}

这说明 $M(\lambda)$ 有唯一的不可约商 $\g$-模 $L(\lambda)$。并且 $L(\lambda)^{\lambda-\rho}$ 是一维的。我们称 $\lambda-\rho$ 为 $L(\lambda)$ 的\emph{最高权}。

\begin{proposition}\label{prop:05.1.6}
\begin{enumerate}[label=(\roman*)]
\item 任一不可约最高权模都同构于某个 $L(\lambda)$。
\item $L(\lambda)\cong L(\mu)$ 当且仅当 $\lambda=\mu$。
\end{enumerate}
\end{proposition}

\begin{proof}
(i) 设 $V$ 是不可约最高权模，$S$ 是 $V$ 的全部权构成的集合。由 1.1，可取权 $\lambda\in S$，使对任意 $\alpha\in\Sigma^+$ 都有 $\lambda+\alpha\notin S$。这说明 $V^\lambda$ 被 $\n_0$ 零化。于是 $V^\lambda$ 是 $\bb_0$-不变的，且其中含有一个 $\bb_0$-不变的一维子空间。设 $v$ 是该子空间中的非零向量。则同态 $U(\g)\to V$，$\xi\mapsto\xi\cdot v$，是满射，且通过 $M(\lambda+\rho)$ 分解。这说明 $V\cong L(\lambda+\rho)$。

(ii) 由 1.5.(i) 与 (ii) 即得。
\end{proof}

\begin{lemma}\label{lem:05.1.7}
中心 $Z(\g)$ 通过特征标 $\chi_\lambda$ 作用在 $M(\lambda)$ 上。
\end{lemma}

\begin{proof}
由 Harish-Chandra 同态的定义立见。
\end{proof}

\begin{proposition}\label{prop:05.1.8}
最高权模有有限长度。
\end{proposition}

\begin{proof}
设 $V$ 是最高权模。若 $V$ 不是有限长度，则可构造 $V$ 的递降 $\g$-模滤过 $(V_i;\,i\in\Z_+)$，使各商 $V_i/V_{i+1}$（$i\in\Z_+$）都不可约。于是由 1.6.(i)，$V_i/V_{i+1}=L(\lambda_i)$，其中 $\lambda_i\in\h_0^*$。由 1.3 与 1.7，可能的 $\lambda_i$ 只有有限多个。再由 1.6.(ii)，可能的 $L(\lambda_i)$ 也只有有限多个，这与 $V$ 的权子空间有限维相矛盾。
\end{proof}

由 [5, Ch. VIII, \S 5, Prop. 2]，存在 $\g$ 的对合自同构 $\iota$，使 $\iota|_{\h_0}=-1$。于是对任意 $\alpha\in R_0^+$ 有 $\iota(\g_\alpha)=\g_{-\alpha}$。设 $\tau$ 是 $U(\g)$ 的反自同构，它是 $U(\g)$ 的主反自同构与「延拓 $\iota$ 的自同构」之积。则 $\tau$ 在 $\h_0$ 上是恒等，且把 $\n_0$ 映入 $\bar{\n}_0$。

\begin{lemma}\label{lem:05.1.9}
反自同构 $\tau$ 在 $Z(\g)$ 上是恒等。
\end{lemma}

\begin{proof}
设 $\gamma$ 是 Harish-Chandra 同态，即 $Z(\g)\subset U(\h_0)\oplus\n_0U(\g)$ 沿 $\n_0U(\g)$ 到 $U(\h_0)$ 的投影。按定义，反自同构 $\tau$ 在 $U(\h_0)$ 上是恒等，且把 $\n_0U(\g)$ 映入 $U(\g)\bar{\n}_0$。另一方面，$\n_0U(\g)$ 与 $U(\g)\bar{\n}_0$ 同 $U(\g)$ 中 $\h_0$ 的中心化子的交相等；这立刻推出 $\gamma$ 与 $\gamma\circ\tau$ 在 $Z(\g)$ 上一致。由 $\gamma$ 的单性即知 $\tau|_{Z(\g)}$ 是恒等。
\end{proof}

设 $V$ 是最高权模，$V^*$ 是它的线性对偶。通过
$$(\xi f)(v)=f(\tau(\xi)v),\qquad \xi\in U(\g),\ f\in V^*,\ v\in V,$$
定义 $U(\g)$ 在 $V^*$ 上的作用。这样 $V^*$ 成为 $\g$-模。设 $\tilde V$ 是由一切使 $U(\h_0)f$ 有限维的 $f\in V^*$ 构成的子空间。容易验证 $\tilde V$ 是 $\g$-子模。

\begin{lemma}\label{lem:05.1.10}
\begin{enumerate}[label=(\roman*)]
\item $\tilde V$ 是最高权模；
\item $(\tilde V)^{\sim}=V$；
\item 对任意 $\lambda\in\h_0^*$ 有 $(\tilde V)^\lambda=(V^\lambda)^*$。
\end{enumerate}
\end{lemma}

\begin{proof}
显然，作为 $\h_0$-模有 $\tilde V=\bigoplus_\lambda(V^\lambda)^*$，且 $(\tilde V)^\lambda=(V^\lambda)^*$。故 $\tilde V$ 的权集与 $V$ 的权集相同。从 $V$ 到 $(\tilde V)^{\sim}$ 的典范映射是单射，也是 $\g$-模态射，且对任意 $\lambda\in\h_0^*$ 有
$$\dim\bigl((\tilde V)^{\sim}\bigr)^\lambda=\dim(\tilde V)^\lambda=\dim V^\lambda,$$
故它是同构。设 $U$ 是 $\tilde V$ 的 $\g$-子模，$U^\perp$ 是 $(\tilde V)^{\sim}=V$ 中与 $U$ 正交的子空间。则 $U^\perp$ 是 $V$ 的 $\g$-子模，且 $(U^\perp)^\perp=U$。这说明 $\tilde V$ 的每个 $\g$-子模都是 $V$ 的某个 $\g$-子模的正交补。由 1.8 可知 $\tilde V$ 有有限长度。特别地，$\tilde V$ 有限生成，从而由 1.1 是最高权模。
\end{proof}

于是 $V\mapsto\tilde V$ 是从 $\mathcal{M}_{fg}(\g,N_0)$ 到自身的正合反变函子。我们称 $\tilde V$ 为 $V$ 的\emph{对偶}。此外，对 Weyl 群 $W$ 在 $\h^*$ 中的任一轨道 $\theta$，由 1.9 可知：$V\in\mathcal{M}_{fg}(U_\theta,N_0)$ 蕴含 $\tilde V\in\mathcal{M}_{fg}(U_\theta,N_0)$。也就是说，$V\mapsto\tilde V$ 是范畴 $\mathcal{M}_{fg}(U_\theta,N_0)$ 与自身的反等价。

\begin{lemma}\label{lem:05.1.11}
对任意 $\lambda\in\h_0^*$ 有 $\widetilde{L(\lambda)}=L(\lambda)$。
\end{lemma}

\begin{proof}
由 1.10，$\widetilde{L(\lambda)}$ 是不可约最高权模，最高权为 $\lambda$。由 1.6 即知 $\widetilde{L(\lambda)}\cong L(\lambda)$。
\end{proof}

今置 $I(\lambda)=\widetilde{M(\lambda)}$。则由 1.5 与 1.11，$I(\lambda)$ 有唯一的不可约 $\g$-子模 $L(\lambda)$。模 $I(\lambda)$ 具有下述泛性质。

\begin{lemma}\label{lem:05.1.12}
设 $\lambda\in\h^*$，$\theta=W\cdot\lambda$。
\begin{enumerate}[label=(\roman*)]
\item 设 $V$ 是最高权模，$\lambda-\rho$ 是 $V$ 的权，且对任意正根 $\alpha\in\Sigma^+$，$\lambda-\rho+\alpha$ 都不是 $V$ 的权。则存在从 $V$ 到 $I(\lambda)$ 的非零态射。
\item 设最高权模 $V$ 满足：
\begin{enumerate}[label=(\alph*)]
\item $V$ 含唯一的不可约子模，它同构于 $L(\lambda)$；
\item 对任意 $\mu\in\h^*$ 有 $\dim V^\mu=\dim I(\lambda)^\mu$。
\end{enumerate}
则 $V\cong I(\lambda)$。
\end{enumerate}
\end{lemma}

\begin{proof}
(i) 由 1.10，$\tilde V$ 是最高权模，$\lambda-\rho$ 是 $\tilde V$ 的权，且对任意正根 $\alpha\in\Sigma^+$，$\lambda-\rho+\alpha$ 都不是 $\tilde V$ 的权。于是若 $v\in(\tilde V)^{\lambda-\rho}$、$v\neq 0$，则 $v$ 被 $\n_0$ 零化。因此同态 $\xi\mapsto\xi\cdot v$ 从 $U(\g)$ 映入 $\tilde V$，且通过 $M(\lambda)$ 分解；于是得到从 $M(\lambda)$ 到 $\tilde V$ 的非零态射 $\varphi$。取对偶，$\tilde\varphi$ 即从 $V$ 到 $I(\lambda)$ 的非零态射。
\begin{annote}
底本此处把目标模 $\tilde V$ 印作 $V$。但 $v$ 取自 $(\tilde V)^{\lambda-\rho}$，且同句后半即写作从 $M(\lambda)$ 到 $\tilde V$ 的态射，故应为 $\tilde V$。已改。[来源：底本 p.97，1.12 证明的上下文自证]
\end{annote}

(ii) 由 (i)，存在从 $V$ 到 $I(\lambda)$ 的非零态射 $\varphi$。因其像非零，必含 $I(\lambda)$ 的唯一不可约子模 $L(\lambda)$。记 $K=\ker\varphi$。由 (b) 有 $\dim V^{\lambda-\rho}=\dim I(\lambda)^{\lambda-\rho}=1$，又 $\dim L(\lambda)^{\lambda-\rho}=1$，故 $K^{\lambda-\rho}=0$。由 (a) 即知 $K=0$，从而 $\varphi$ 是单射。最后由 (b) 可知 $\varphi$ 是同构。
\end{proof}

现在把上述结果与 $\mathcal{M}_{fg}(U_\theta,N_0)$ 中不可约对象分类的一般几何方案联系起来。首先，由 Bruhat 引理（[2], 14.11），$N_0$ 在旗簇 $X$ 中只有有限多个轨道。因此有下述注记；借助它，可以应用……的结果。

\begin{proposition}\label{prop:05.1.13}
$(\g,N_0)$ 是 Harish-Chandra 对。
\end{proposition}

$X$ 中的 $N_0$-轨道就是 Bruhat 胞腔 $C(w)$（$w\in W$）。它们是 $X$ 的仿射子簇（[2], 14.11），且
$$\dim C(w)=\ell(w),\qquad w\in W.$$
因此，若 $C(s)$ 落在 $C(w)$ 的边界 $\partial C(w)=\overline{C(w)}-C(w)$ 中，则 $\ell(s)<\ell(w)$。

设 $\lambda\in\h^*$。设 $C(w)$ 是 Bruhat 胞腔，$i_w:C(w)\to X$ 是典范浸入。则由……可见，$C(w)$ 上唯一的不可约 $N_0$-齐性 $(\D_\lambda)_{i_w}$-联络是 $\calO_{C(w)}$。与数据 $(C(w),\calO_{C(w)})$ 相应的标准 $\D_\lambda$-模记作 $I(w,\lambda)$，它唯一的不可约 $\D_\lambda$-子模记作 $L(w,\lambda)$。

不可约对象的几何分类与 1.6 之间的关键联系，由下述结果给出。
\begin{annote}
底本此处印作「and 6.」，按上下文应为 1.6（命题 \ref{prop:05.1.6}）：它给出不可约最高权模的代数分类，正是本段所说的几何分类在代数侧的对应物。已按 1.6 处理。[来源：底本 p.97；上下文指涉]
\end{annote}

\begin{theorem}\label{thm:05.1.14}
设 $\lambda\in\h^*$ 反支配。则
$$\Gamma(X,I(w,\lambda))=I(w\lambda),\qquad w\in W.$$
\end{theorem}

为证 1.14，先作些准备。我们从 1.14 的一个非常特殊的情形开始。

\begin{lemma}\label{lem:05.1.15}
设 $\lambda\in\h^*$ 反支配。则
$$\Gamma(X,I(1,\lambda))=M(\lambda)=L(\lambda)=I(\lambda).$$
\end{lemma}

\begin{proof}
显然 $I(1,\lambda)$ 是不可约 $\D_\lambda$-模。故由 L.4.1，$U_\theta$-模 $\Gamma(X,I(1,\lambda))$ 或者不可约最高权模，或者为零。由于 $I(1,\lambda)$ 的支集是点 $C(1)$，后一可能自动排除。因此，为证本命题，只需建立第一个等号。

下面描述直像模 $I(1,\lambda)=i_{1+}(\C_{\lambda+\rho})$ 的结构。可把它看作右 $\D_{-\lambda}$-模。此时它按自然方式等于右 $\D_{-\lambda}$-模 $i_1^*(\D_{-\lambda,C(1)\to X})$。另一方面，如同 L.2.4 的证明，可得作为右 $U(\g)$-模
\begin{align*}
\D_{-\lambda,C(1)\to X}&=T_{x_0}(\D_{-\lambda})\\
&=(U(\g)/\n_0U(\g))\big/\bigl(I_{-\lambda+\rho}(U(\g)/\n_0U(\g))\bigr)\\
&=\C_{-\lambda+\rho}\otimes_{U(\bb_0)}U(\g),
\end{align*}
其中 $x_0$ 记 $C(1)$ 中的点。这说明作为左 $U(\g)$-模，$I(1,\lambda)$ 等于 $M(\lambda)$。
\end{proof}

下面需要交织函子作用在标准模\index{biaozhunmo@标准模|hyperpage}上的一些结果。

\begin{lemma}\label{lem:05.1.16}
设 $w\in W$，$\lambda\in\h^*$。则
$$LI_w(D(I(1,\lambda)))=D(I(w^{-1},w\lambda)).$$
\end{lemma}

\begin{proof}
沿用 L.3 的记号。设 $Z_w\subset X\times X$ 是处于相对位置 $w\in W$ 的有序 Borel 子代数对构成的簇。记 $p_i$（$i=1,2$）为 $X\times X$ 到第 $i$ 个因子的投影。则
\begin{align*}
p_2^{-1}(C(1))&=\{(x,x')\in X\times X\mid \bb_{x'}=\bb_0,\ \bb_x\ \text{相对于}\ \bb_0\ \text{处于相对位置}\ w^{-1}\}\\
&=C(w^{-1})\times C(1),
\end{align*}
即有交换图表
\begin{center}
\begin{tikzcd}
C(w^{-1})\times C(1) \arrow[r,"j"] \arrow[d,"pr_2"'] & Z_w \arrow[d,"p_2"]\\
C(1) \arrow[r,"i_1"'] & X
\end{tikzcd}
\end{center}
由基变换（[1], VI.8.4），因 $pr_2$ 与 $p_2$ 是淹没、$i_1$ 与 $j$ 是仿射浸入，有
$$p_2^+(I(1,\lambda))=p_2^+(i_{1+}(\calO_{C(1)}))=j_+(pr_2^+(\calO_{C(1)}))=j_+(\calO_{C(w^{-1})\times C(1)}).$$
投影 $p_1$ 诱导 $p_2^{-1}(C(1))$ 到 $X$ 的浸入，其像等于 $C(w^{-1})$，即有交换图表
\begin{center}
\begin{tikzcd}
C(w^{-1})\times C(1) \arrow[r,"j"] \arrow[d,"pr_1"'] & Z_w \arrow[d,"p_1"]\\
C(w^{-1}) \arrow[r,"i_{w^{-1}}"'] & X
\end{tikzcd}
\end{center}
核对相应的扭之后得到
\begin{align*}
LI_w(D(I(1,\lambda)))&=Rp_{1+}(T_w\otimes_{\calO_{Z_w}}p_2^+(I(1,\lambda)))\\
&=Rp_{1+}(T_w\otimes_{\calO_{Z_w}}j_+(\calO_{C(w^{-1})\times C(1)}))\\
&=D(i_{w^{-1}+}(\calO_{C(w^{-1})}))=D(I(w^{-1},w\lambda)).
\end{align*}
\end{proof}

\begin{corollary}\label{cor:05.1.17}
设 $w,w'\in W$ 满足 $\ell(ww')=\ell(w)+\ell(w')$。则
$$LI_w(D(I(w'^{-1},\lambda)))=D(I(w'^{-1}w^{-1},w\lambda)).$$
\end{corollary}
